\section{一个 bit 只回答两种状态}

先看一只开关。它可以断开，也可以闭合；如果我们只关心这两种可区分状态，
就可以把它们分别记作 0 和 1。这里的 0 和 1 是状态的名字，不要求开关内部
真的存在“数字”。计算机中的晶体管、电压和磁性区域都可以承担这两个状态，
后面几章再讨论它们怎样做到。

\term{bit}{binary digit，二进制位}是只能取 0 或 1 的最小状态位置。一个 bit
有两种取值，两个互相独立的 bit 有 \(2\times2=2^2\) 种组合，\(n\) 个 bit
共有
\[
  2^n
\]
种组合。这个公式来自逐位选择：每增加一个 bit，组合数就再乘二。

\begin{osgridtabular}{L{0.16\textwidth}L{0.24\textwidth}L{0.45\textwidth}}
bit 数 & 组合数 & 可以编号的范围\\ \hline
1 & \(2\) & 0--1\\ \hline
4 & \(16\) & 0--15\\ \hline
8 & \(256\) & 0--255\\ \hline
16 & \(65536\) & 0--65535\\ \hline
\end{osgridtabular}

\section{八个 bit 组成一个 byte}

现代 x86 把八个连续 bit 作为一个 \term{byte}{字节}。一个字节有
\(2^8=256\) 种组合，可以表示无符号整数 0--255。存储容量中的 B 表示
byte，b 表示 bit；因此 \(1\,\mathrm{B}=8\,\mathrm{b}\)。网络速率写成
Mbit/s 时，不能直接把数值当成 MB/s。

二进制数 \(11001010_2\) 的每一位都有权重：
\[
  1\cdot2^7+1\cdot2^6+0\cdot2^5+0\cdot2^4+
  1\cdot2^3+0\cdot2^2+1\cdot2^1+0\cdot2^0=202.
\]
手算一次比背诵转换表更重要，因为寄存器位、页表项和设备状态最终都要按同样
方法拆开。

\section{十六进制让字节更容易阅读}

四个 bit 恰好有 16 种组合，所以可以用一个十六进制数字表示。二进制
\(1100\,1010_2\) 分成两组后得到 \(C\,A\)，写成 \code{0xCA}。一个字节
总能写成两个十六进制数字，两个字节总能写成四个。

\begin{osgridtabular}{L{0.18\textwidth}L{0.18\textwidth}L{0.36\textwidth}}
二进制 & 十六进制 & 十进制\\ \hline
\code{0000} & \code{0} & 0\\ \hline
\code{1001} & \code{9} & 9\\ \hline
\code{1010} & \code{A} & 10\\ \hline
\code{1111} & \code{F} & 15\\ \hline
\end{osgridtabular}

十六进制没有改变数据。它只缩短了人在纸面上书写同一组 bit 的方式。后面看到
\code{0x3F8}、\code{0xFFFFFFF0} 或 \code{0x20} 时，都可以随时展开回二进制
检查其中哪一位为 1。

\section{地址给每个字节一个编号}

内存可以先想成一排带编号的小格，每格保存一个字节。格子的编号叫地址。
若地址从 0 开始，容量为 \(N\) 字节，那么有效地址是
\[
  0,\ 1,\ \ldots,\ N-1.
\]
最后一个地址不是 \(N\)，因为 0 已经占用了第一个编号。

连续范围写成半开区间 \([begin,end)\)。起点属于范围，终点不属于范围，
字节数直接等于 \(end-begin\)。例如 \([0x1000,0x1800)\) 包含
\[
  0x1800-0x1000=0x800=2048
\]
字节。这个写法以后会反复用于 ROM、栈、ELF 段和虚拟内存区域。

\section{同一组字节可以有不同含义}

字节 \code{0x41} 可以解释成无符号整数 65，也可以按 ASCII 解释成字符 A；
放在 x86 指令的特定位置时，它还可能成为立即数或地址的一部分。CPU 不会替
我们猜测含义。含义来自读取这些字节的规则。

本章结束时只需要能够完成三件事：把二进制按四位分组写成十六进制；区分 bit
与 byte；由起止地址算出容量。下一章开始解释真实电路怎样产生可以稳定区分的
0 和 1。
